\documentclass{article}


\usepackage[final]{neurips_2024}

\usepackage{subcaption}
\usepackage[utf8]{inputenc} % allow utf-8 input
\usepackage[T1]{fontenc}    % use 8-bit T1 fonts
\usepackage{hyperref}       % hyperlinks
\usepackage{url}            % simple URL typesetting
\usepackage{booktabs}       % professional-quality tables
\usepackage{amsfonts}       % blackboard math symbols
\usepackage{nicefrac}       % compact symbols for 1/2, etc.
\usepackage{microtype}      % microtypography
\usepackage{xcolor}         % colors
\usepackage{amsmath}        % for \eqref and math environments
\usepackage{multirow}      % multi-row cells in tables
\usepackage{graphicx}
\usepackage{array}
\usepackage{longtable}

% \usepackage[nonatbib]{neurips_2024}
% \newcommand*\samethanks[1][\value{footnote}]{\footnotemark[#1]}
\title{VisText-VLM2Plastic: Evaluating Multilingual and Image-Grounded Environmental Reasoning in Vision-Language Models}

% \author{%
% author $^{2}$\thanks{Equal Contribution} \quad author$^{1,3*}$ \quad author$^{3}$ \\ 
% $^1$United International University \quad $^2$American International University Bangladesh \quad $^3$lAB NAME\\
% \texttt{m1}
% \texttt{m2}
% \texttt{m3}
% }


\begin{document}


\maketitle

\vspace{-0.6cm}

\begin{abstract}
Vision-language models (VLMs) increasingly support open-ended visual reasoning, yet their behavior in low-resource languages and sustainability-oriented domains remains underexplored. We introduce a controlled English--Bangla plastic-waste VQA benchmark with 402 image records and 2,412 question-answer pairs across five waste categories and three reasoning levels: recognition, visual description, and environmental reasoning. Four VLMs GPT-4.1-mini, Gemini 2.5 Flash, Qwen2.5-VL-7B-Instruct, and Gemma-3-4B are evaluated zero-shot using recognition, generation, bilingual-consistency, and visual-ablation metrics. Gemini 2.5 Flash achieves the best recognition performance with 93.5\% accuracy and 91.3\% macro-F1. Across models, English performance exceeds Bangla by an average of 3.78 points, with the largest gap in environmental reasoning. Visual-ablation gains are 66.29 points for recognition, 52.89 for description, and only 11.42 for environmental reasoning, indicating substantially weaker image dependence for the latter. These results show that conventional answer-quality metrics can mask weaknesses in multilingual robustness and visual grounding, motivating evaluation frameworks that jointly measure answer quality, cross-lingual consistency, and dependence on visual evidence.
\end{abstract}

\section{Introduction}
\label{sec:introduction}

Plastic waste creates persistent environmental burdens through litter accumulation, drainage blockage, and movement into waterways. Computer vision has therefore been widely explored for waste classification and sorting, while recent vision-language models (VLMs) provide a more flexible alternative by jointly processing images and natural-language instructions. \citet{malla2025wasterecognition} showed that contemporary VLMs can achieve competitive zero-shot waste recognition and benefit substantially from prompt engineering.

Recognition alone, however, does not establish environmental understanding. A useful system should be able to identify the visible waste item, describe only visually supported properties, and connect the observed object to scientifically reasonable environmental concerns. These abilities are distinct: a model may correctly recognize an object while hallucinating unsupported visual details or producing generic environmental explanations that do not depend on the image.

Visual Question Answering (VQA) provides a natural framework for studying these distinctions \citep{antol2015vqa}. However, prior work has shown that high VQA accuracy can arise from language priors rather than genuine visual understanding \citep{goyal2017makingv}. This issue is particularly relevant for environmental questions, which may often be answered plausibly from pretrained knowledge without using the supplied image.

The problem becomes more important in multilingual settings. Bangla remains underrepresented in multimodal evaluation despite recent resources such as ChitroJera \citep{barua2024chitrojera}, BanglaProtha \citep{fahim2026banglaprotha}, and BanglaMedVQA \citep{ahmed2026banglamedvqa}. These studies show that strong general-purpose multimodal systems can degrade in specialized Bangla contexts.

We therefore develop a controlled English--Bangla benchmark for plastic-waste understanding organized into three reasoning levels:
\begin{equation}
\text{Recognition}
\;\longrightarrow\;
\text{Visual Description}
\;\longrightarrow\;
\text{Environmental Reasoning}.
\label{eq:reasoning_hierarchy}
\end{equation}

The benchmark addresses three research questions:
\begin{itemize} \item[\textbf{RQ1.}] How does VLM performance change as the task progresses from direct plastic-waste recognition to visual description and environmental reasoning? \item[\textbf{RQ2.}] How large is the performance gap between semantically aligned English and Bangla questions, and does this gap increase with reasoning complexity? \item[\textbf{RQ3.}] To what extent are apparently correct environmental answers dependent on the image rather than recoverable from language priors alone? \end{itemize}

Our contributions are:
\begin{itemize}
    \item We introduce a controlled English--Bangla diagnostic benchmark centered on plastic-waste understanding rather than general-purpose or culturally oriented VQA, extending recent work on Bangla multimodal evaluation.

    \item Evaluate proprietary and open VLMs across recognition, visual description, and environmental reasoning using classification, lexical, semantic, and cross-lingual metrics.

    \item Introduce an explicit image-ablation protocol that quantifies how much each task benefits from visual information and distinguishes visually grounded reasoning from language-prior-driven responses, following the broader motivation of visually grounded VQA evaluation.
    \item We analyze bilingual consistency and model failure modes, demonstrating why multilingual environmental VQA cannot be adequately characterized using a single aggregate performance score.
\end{itemize}

\section{Related Work}
\label{sec:related_work}

\subsection{Visual Question Answering and Language Priors}

VQA requires models to jointly reason over images and natural-language questions \citep{antol2015vqa}. However, \citet{goyal2017makingv} showed that models can exploit question regularities and language priors, producing strong scores without sufficiently interpreting the image. Our work follows this diagnostic perspective but directly removes the image to quantify modality dependence across reasoning levels.

\subsection{Bangla Multimodal Evaluation}

Bangla multimodal evaluation remains comparatively limited. ChitroJera introduced regionally relevant Bangla VQA data \citep{barua2024chitrojera}, BanglaProtha studied Bengali cultural understanding across VLMs \citep{fahim2026banglaprotha}, and BanglaMedVQA evaluated domain-specific medical reasoning in Bangla \citep{ahmed2026banglamedvqa}. Our benchmark differs by focusing on environmental reasoning and by using semantically aligned English--Bangla questions for the same images.

\subsection{Vision-Language Models for Waste Understanding}

Recent work has extended waste analysis beyond supervised image classification. \citet{malla2025wasterecognition} evaluated VLMs for zero-shot waste recognition and reported substantial gains from prompt optimization. WasteAssistant further introduced waste-oriented VQA and regulation-aware reasoning across multiple waste categories \citep{kataruka2026wasteassistant}. In contrast, our benchmark focuses on multilingual robustness and visual dependence, asking whether Bangla performance differs from English and whether environmental responses genuinely require image evidence.
\begin{table}[t] \centering \caption{Benchmark composition by plastic-waste category.} \label{tab:benchmark_composition} \begin{tabular}{lrr} \toprule \textbf{Category} & \textbf{Image Records} & \textbf{Fraction} \\ \midrule Plastic wrapper / flexible packaging & 224 & 55.7\% \\ Plastic bottle & 70 & 17.4\% \\ Plastic cup & 48 & 11.9\% \\ Plastic straw & 36 & 9.0\% \\ Plastic bag / plastic sheet & 24 & 6.0\% \\ \midrule \textbf{Total} & \textbf{402} & \textbf{100\%} \\ \bottomrule \end{tabular} \end{table}
\section{Dataset Benchmark}
\label{sec:benchmark}
\begin{figure}[t]
    \centering
    \includegraphics[width=\linewidth]{figures/benchmark_pipeline.png}
    \caption{Benchmark and evaluation pipeline. Each grouped plastic-waste image produces aligned English and Bangla recognition, description, and environmental-reasoning questions. Four VLMs are evaluated under image+question and question-only conditions, followed by multilingual, grounding, and error analyses.}
    \label{fig:benchmark_pipeline}
\end{figure}
The benchmark contains 402 image records derived from 136 underlying base images. Multiple records correspond to augmented variants of the same source image; therefore, a \texttt{source\_base\_id} is preserved for every record, and all augmentations originating from the same source image are treated as a single group during dataset splitting. The final benchmark contains five dominant plastic-waste categories, whose distribution is summarized in Table~\ref{tab:benchmark_composition}. The dataset is imbalanced, with flexible wrappers comprising more than half of the image records, motivating the use of macro-averaged metrics in addition to micro accuracy. The dataset is split into 282/58/62 train, validation, and test image records, corresponding to 95/20/21 base images and 1,692/348/372 QA pairs, respectively, and no \texttt{source\_base\_id} appears in more than one partition. This grouping prevents visually near-identical augmented samples from appearing across development and evaluation sets, which could otherwise inflate performance. The complete benchmark and evaluation flow is summarized in Figure~\ref{fig:benchmark_pipeline}.

Each image record is associated with six QA pairs: English recognition, Bangla recognition, English visual description, Bangla visual description, English environmental reasoning, and Bangla environmental reasoning. This produces 2,412 QA pairs, evenly divided into 1,206 English and 1,206 Bangla instances. Question construction follows a controlled template-based procedure in which surface wording varies while task semantics are preserved. All 2,412 questions have unique surface strings in the current release. Reference answers are generated from a constrained category-conditioned answer bank containing 100 distinct answer strings across the two languages. This controlled design reduces uncontrolled stylistic variation while limiting linguistic diversity; this trade-off is discussed further in Appendix~\ref{sec:limitations}. The overall data-collection and construction workflow is illustrated in Appendix \ref{appendix:dataset} Figure~\ref{fig:data_collection_pipeline}.

% \begin{table}[t]
% \centering
% \caption{Dataset split statistics.}
% \label{tab:dataset_split}
% \begin{tabular}{lrrr}
% \toprule
% \textbf{Split} & \textbf{Image Records} & \textbf{Base Images} & \textbf{QA Pairs} \\
% \midrule
% Train      & 282 & 95  & 1,692 \\
% Validation & 58  & 20  & 348 \\
% Test       & 62  & 21  & 372 \\
% \midrule
% \textbf{Total} & \textbf{402} & \textbf{136} & \textbf{2,412} \\
% \bottomrule
% \end{tabular}
% \end{table}

\begin{table}[httb]
\centering
\caption{Zero-shot recognition performance.}
\label{tab:recognition_results}
\begin{tabular}{lcc}
\toprule
\textbf{Model} & \textbf{Accuracy (\%)} & \textbf{Macro-F1 (\%)} \\
\midrule
GPT-4.1 mini
& 91.9
& 89.2 \\

Gemini 2.5 Flash
& \textbf{93.5}
& \textbf{91.3} \\

Qwen2.5-VL-7B
& 88.7
& 85.6 \\

Gemma 3 4B
& 85.5
& 81.4 \\
\bottomrule
\end{tabular}
\end{table}
\section{Experimental Setup}
\label{sec:experimental_setup}
Experiments were conducted in Google Colab using an NVIDIA Tesla T4 GPU, with the pipeline implemented in Python using the required deep-learning and vision-language libraries. The T4 GPU was used for dataset preprocessing, local inference with the open-source VLMs, and metric computation, while the proprietary models were accessed through their respective APIs. We evaluate four representative VLMs spanning proprietary and openly available model families: GPT-4.1 mini, using the fixed API snapshot \texttt{gpt-4.1-mini-2025-04-14} where available, which supports image inputs and improved multimodal understanding and instruction following \citep{openai2025gpt41}; Gemini 2.5 Flash, which serves as a proprietary multimodal baseline optimized for efficient reasoning \citep{google2025gemini25flash}; Qwen2.5-VL-7B-Instruct, which supports dynamic-resolution visual processing and multilingual vision-language interaction \citep{bai2025qwen25vl}; and the multimodal Gemma-3-4B Instruction-Tuned model, which provides vision understanding and expanded multilingual capabilities \citep{gemmateam2025gemma3}. The comparison is intended to represent different deployment regimes, model scales, and accessibility conditions rather than to establish a universal ranking, since model families and API implementations may evolve over time. Additional details regarding the experimental setup, and evaluation procedures are provided in Appendix~\ref{appendix:metrics}.


\section{Results}
\label{sec:results}


\subsection{Zero-Shot Recognition}
\label{subsec:recognition_results}

Table~\ref{tab:recognition_results} reports the zero-shot recognition performance of all four evaluated models.
The experimental results suggest that recognition is the strongest task across all four models. Gemini 2.5 Flash achieves the highest category accuracy at 93.5\%, followed by GPT-4.1 mini at 91.9\%. The difference between accuracy and macro-F1 is consistent with the class imbalance described in Section~\ref{sec:benchmark}: models benefit disproportionately when they perform well on the dominant flexible-wrapper class.

Category-level results complement the aggregate metrics and reveal greater sensitivity for bag and straw examples, which have fewer samples and can visually overlap with flexible packaging. Figure~\ref{fig:category_performance} presents the category distribution together with per-category recognition performance.
\begin{table*}[t]
\centering
\caption{Open-ended response quality across description and environmental-reasoning tasks. }
\label{tab:generation_results}
\begin{tabular}{lcccccc}
\toprule
&
\multicolumn{3}{c}{\textbf{Description}}
&
\multicolumn{3}{c}{\textbf{Environmental Reasoning}}
\\
\cmidrule(lr){2-4}
\cmidrule(lr){5-7}
\textbf{Model}
&
\textbf{BERT-F1}
&
\textbf{BLEU-4}
&
\textbf{ROUGE-L}
&
\textbf{BERT-F1}
&
\textbf{BLEU-4}
&
\textbf{ROUGE-L}
\\
\midrule

GPT-4.1 mini
& 0.868
& 0.431
& 0.619
& 0.840
& 0.347
& 0.566 \\

Gemini 2.5 Flash
& \textbf{0.881}
& \textbf{0.452}
& \textbf{0.641}
& \textbf{0.855}
& \textbf{0.369}
& \textbf{0.588} \\

Qwen2.5-VL-7B
& 0.834
& 0.398
& 0.587
& 0.802
& 0.311
& 0.531 \\

Gemma 3 4B
& 0.804
& 0.361
& 0.552
& 0.769
& 0.276
& 0.498 \\

\bottomrule
\end{tabular}
\end{table*}
\begin{figure}[t]
    \centering
    \includegraphics[width=\linewidth]{figures/category_performance.png}
    \caption{Category distribution and per-category recognition performance. Error bars represent 95\% cluster-bootstrap confidence intervals over source base images.}
    \label{fig:category_performance}
\end{figure}

\subsection{Description and Environmental Generation}
\label{subsec:generation_results}
Table~\ref{tab:generation_results} summarizes BERTScore-F1, BLEU-4, and ROUGE-L for the description and environmental-reasoning tasks. The observed pattern is consistent across all three automatic metrics: visual description is easier than environmental reasoning, and the two proprietary systems outperform the selected 7B and 4B open checkpoints. BERTScore values are substantially higher than BLEU-4 scores. This is expected rather than contradictory. A model may describe a wrapper as ``thin flexible packaging'' when the reference uses ``lightweight plastic wrapper with a sheet-like form.'' The two responses preserve similar semantic content while sharing relatively few four-grams. This distinction motivates using BERTScore as the primary open-ended similarity measure \citep{zhang2020bertscore}, with BLEU-4 and ROUGE-L serving as complementary lexical metrics \citep{papineni2002bleu,lin2004rouge}.

The reduction from description to environmental reasoning is especially important. Environmental answers require not only recognizing the visible object but also connecting it with appropriate domain knowledge while avoiding unsupported toxicological, chemical, or causal claims.
\begin{figure}[t]
    \centering
    \includegraphics[width=\linewidth]{figures/language_gap.png}
    \caption{English--Bangla performance gap across recognition, description, and environmental reasoning. Error bars represent 95\% cluster-bootstrap confidence intervals over source base images.}
    \label{fig:language_gap}
\end{figure}

\begin{figure}[t]
    \centering
    \includegraphics[width=\linewidth]{figures/grounding_ablation.png}
    \caption{Visual-ablation analysis. Recognition and description exhibit substantially larger multimodal gains than environmental reasoning in the experimental results.}
    \label{fig:grounding_ablation}
\end{figure}
\subsection{English--Bangla Performance Gap}
\label{subsec:language_gap}
We next examine whether semantically aligned English and Bangla questions produce equivalent model behavior. Model-level normalized performance and the corresponding English--Bangla gap are reported in Table~\ref{tab:language_model_gap}.
\begin{table}[t]
\centering
\caption{Language-level normalized performance. }
\label{tab:language_model_gap}
\begin{tabular}{lccc}
\toprule
\textbf{Model}
& \textbf{English}
& \textbf{Bangla}
& \textbf{EN--BN Gap} \\
\midrule
GPT-4.1 mini
& 89.23
& 85.87
& 3.36 \\

Gemini 2.5 Flash
& \textbf{90.60}
& \textbf{87.43}
& \textbf{3.17} \\

Qwen2.5-VL-7B
& 85.90
& 82.23
& 3.67 \\

Gemma 3 4B
& 83.10
& 78.20
& 4.90 \\

\midrule
\textbf{Mean}
& --
& --
& \textbf{3.78} \\
\bottomrule
\end{tabular}
\end{table}
All four evaluated models exhibit an English advantage. The gap is smallest for Gemini 2.5 Flash and largest for Gemma 3 4B. More informative than the overall difference is how the language gap changes with reasoning complexity; Table~\ref{tab:language_task_gap} reports this comparison by task.
\begin{table}[t]
\centering
\caption{Language gap by reasoning level. }
\label{tab:language_task_gap}
\begin{tabular}{lccc}
\toprule
\textbf{Task}
& \textbf{English}
& \textbf{Bangla}
& \textbf{Gap} \\
\midrule
Recognition
& 91.53
& 87.90
& 3.63 \\

Description
& 86.43
& 82.88
& 3.55 \\

Environmental reasoning
& 83.68
& 79.53
& \textbf{4.15} \\
\bottomrule
\end{tabular}
\end{table}

Environmental reasoning exhibits the largest language gap, supporting the conclusion that multilingual degradation becomes more visible when models must combine visual understanding with longer-form domain reasoning instead of producing a short category label. Language-specific open-ended scores are reported in Table~\ref{tab:language_generation_metrics}; lexical metrics require additional caution.

\begin{table}[t]
\centering
\caption{Language-specific open-ended evaluation. }
\label{tab:language_generation_metrics}
\begin{tabular}{lccc}
\toprule
\textbf{Language}
& \textbf{BERT-F1}
& \textbf{BLEU-4}
& \textbf{ROUGE-L} \\
\midrule
English
& 0.842
& 0.382
& 0.574 \\

Bangla
& 0.806
& 0.329
& 0.531 \\
\bottomrule
\end{tabular}
\end{table}

The BLEU gap is proportionally larger than the BERTScore gap. This suggests that some Bangla responses may remain semantically appropriate despite greater surface-form variation. BLEU should therefore not be used alone to infer weaker Bangla reasoning. Figure~\ref{fig:language_gap} visualizes the English--Bangla performance gap across the three reasoning levels.




\subsection{Does the Image Actually Matter?}
\label{subsec:image_ablation_results}
Table~\ref{tab:visual_ablation} summarizes the central diagnostic experiment.
\begin{table}[t]
\centering
\caption{Visual-ablation analysis. }
\label{tab:visual_ablation}
\begin{tabular}{lccc}
\toprule
\textbf{Task}
&
\textbf{Image + Q}
&
\textbf{Q Only}
&
\textbf{VGG} \\
\midrule
Recognition
& 89.71
& 23.42
& \textbf{+66.29} \\

Description
& 84.65
& 31.76
& \textbf{+52.89} \\

Environmental reasoning
& 81.60
& 70.18
& \textbf{+11.42} \\
\bottomrule
\end{tabular}
\end{table}
Recognition behaves as expected: removing the image produces a dramatic performance collapse. Description exhibits a similarly large reduction, indicating that the corresponding answers are strongly tied to visible evidence.

Environmental reasoning behaves differently. Even without an image, the observed results retain 70.18 normalized performance compared with 81.60 under multimodal input, yielding a VGG of only 11.42 points. This is a central empirical finding of the evaluation. Environmental questions can contain enough semantic information---for example, ``plastic waste,'' ``improper disposal,'' and ``environmental problem'' to trigger broadly correct statements concerning litter, drainage systems, waterways, or wildlife. Consequently, a high environmental-reasoning score does not, by itself, establish that a model identified or meaningfully reasoned about the visual object.

This behavior is conceptually related to the language-prior problem identified in prior VQA research \citep{goyal2017makingv}. The issue is not necessarily that the generated environmental answer is false; rather, the \emph{source of evidence} supporting that answer may differ from what a multimodal benchmark is intended to measure. The task-wise multimodal and question-only comparison is visualized in Figure~\ref{fig:grounding_ablation}.



\subsection{Cross-Lingual Consistency}
\label{subsec:bcs_results}
A model may achieve respectable aggregate scores in both languages while producing contradictory answers to paired questions. We therefore evaluate cross-lingual response consistency, with model-wise BCS results reported in Table~\ref{tab:bcs_results}.

The observed pattern indicates a monotonic decline in consistency as reasoning complexity increases. Recognition is highly stable because outputs often reduce to the same short waste category. Description introduces greater linguistic variation, while environmental reasoning provides more opportunities for differences in specificity, factual content, and grounding. This analysis captures a failure mode that independent English and Bangla evaluation can conceal. For example, a model may identify an object as a plastic wrapper in English but use a broader or different waste category in Bangla while still generating environmentally plausible follow-up text. Both responses may obtain reasonable independent semantic scores while remaining mutually inconsistent.
\begin{table}[htb]
\centering
\caption{Bilingual Consistency Score (BCS). }
\label{tab:bcs_results}
\begin{tabular}{lccc}
\toprule
\textbf{Model}
&
\textbf{Rec.}
&
\textbf{Desc.}
&
\textbf{Env.} \\
\midrule
GPT-4.1 mini
& 0.94
& 0.89
& 0.84 \\

Gemini 2.5 Flash
& \textbf{0.96}
& \textbf{0.91}
& \textbf{0.87} \\

Qwen2.5-VL-7B
& 0.91
& 0.85
& 0.79 \\

Gemma 3 4B
& 0.87
& 0.81
& 0.73 \\
\bottomrule
\end{tabular}
\end{table}


%
\section{Discussion}
\label{sec:discussion}
%

\subsection{Task Difficulty and Multilingual Robustness}
\label{subsec:task_multilingual_discussion}

The benchmark is intentionally hierarchical. The empirical results show that recognition is the strongest task, while performance decreases when models are required to describe visual evidence or generate environmentally relevant explanations. Recognition accuracy is therefore a necessary but insufficient measure of environmental multimodal capability, particularly for applications such as environmental education, citizen reporting, and waste-assistance systems.

The results also show that multilingual degradation is task dependent. Short recognition answers require comparatively little language generation, whereas environmental questions require longer semantic processing and domain expression. The English--Bangla gap is largest for environmental reasoning, indicating that multilingual weaknesses become more visible as reasoning complexity increases. This observation is consistent with broader findings from low-resource multimodal evaluation; for example, BanglaMedVQA reports that specialized visual reasoning remains challenging even for strong general-purpose LVLMs \citep{ahmed2026banglamedvqa}. A methodological advantage of our benchmark is that English and Bangla questions are paired for the same image and semantic intent, reducing visual-content confounds in cross-lingual comparison.

\subsection{Visual Grounding, Metric Interpretation, and Model Comparison}
\label{subsec:grounding_metric_models}

The image-ablation experiment exposes a central multimodal evaluation problem. If a model can answer an environmental question nearly as well without seeing the image, the benchmark score is partly measuring environmental language knowledge rather than image-grounded reasoning. General environmental knowledge is useful, but it should not be interpreted as evidence that the model understood the supplied image. Future benchmark extensions could reduce this shortcut through counterfactual visual pairs, where similar question structures are paired with different waste categories and reference answers require category-specific consequences.

This grounding issue also affects the interpretation of automatic metrics. BLEU rewards local $n$-gram precision \citep{papineni2002bleu}, ROUGE-L measures longest-common-subsequence overlap \citep{lin2004rouge}, and BERTScore is more tolerant of semantically correct paraphrases \citep{zhang2020bertscore}. These metrics therefore capture different aspects of response quality, especially in bilingual evaluation where Bangla morphology and phrase variation can reduce lexical overlap. None of them directly establishes visual grounding or scientific correctness, so we treat them as complementary and combine them with modality ablation and targeted human evaluation.

The experimental results place Gemini 2.5 Flash and GPT-4.1 mini ahead of the selected open checkpoints. However, this comparison spans different model sizes, training corpora, architectures, inference infrastructures, and deployment constraints, and should not be interpreted as evidence that proprietary VLMs are inherently superior. Open models remain advantageous for reproducibility, local deployment, privacy, modification, and scientific inspection. Future comparisons should include larger open checkpoints and more closely matched compute and latency budgets.

% 


% % 
% \section{Reproducibility}
% \label{sec:reproducibility}
% % 

% The evaluation pipeline exports raw predictions before metric computation. For each prediction, the released result file should contain the model identifier, base-image identifier, image filename, split, category, language, question type, question, reference answer, generated response, and inference condition.

% The accompanying code release should include:

% \begin{itemize}
%     \item dataset loading and grouping by \texttt{source\_base\_id};
%     \item zero-shot prompts and inference configurations;
%     \item exact recognition-normalization rules;
%     \item BERTScore, BLEU-4, ROUGE-L, and LaBSE implementations;
%     \item image-ablation inference;
%     \item source-group cluster-bootstrap confidence intervals;
%     \item bilingual-consistency analysis;
%     \item category-level evaluation; and
%     \item scripts for regenerating all tables and figures.
% \end{itemize}

% For API-based models, raw responses should be cached to reduce instability caused by repeated calls or subsequent model updates. For local models, package versions, checkpoint revisions, quantization settings, GPU type, inference precision, and decoding parameters should be recorded.


% % 
\section{Conclusion}
\label{sec:conclusion}
% 

We introduce a controlled bilingual framework for evaluating how vision-language models interpret plastic waste in English and Bangla. The benchmark contains 402 image records derived from 136 base images and 2,412 paired QA instances spanning recognition, visual description, and environmental reasoning. The framework extends conventional VQA evaluation in three important ways. First, paired English--Bangla questions isolate language-dependent degradation while preserving the underlying visual content and semantic intent. Second, the Bilingual Consistency Score measures whether a model maintains equivalent interpretations across languages. Third, an image-ablation experiment directly evaluates whether environmental responses depend on visual evidence.

The experimental results illustrate a consistent pattern of strong recognition, weaker environmental reasoning, a systematic English--Bangla performance gap, and substantially lower image dependence for environmental reasoning than for recognition or visual description. These findings indicate that contemporary VLMs can possess substantial environmental knowledge without consistently grounding that knowledge in what they see.

This distinction is central to responsible environmental AI. A robust multimodal system should not merely generate plausible statements about pollution. It should identify relevant visual evidence, preserve semantic meaning across languages, avoid unsupported claims, and measurably depend on the image whenever a task requires visual reasoning.
\bibliographystyle{plainnat}
\bibliography{refs}
\clearpage
\appendix
\section{Extra Dataset Information}
\label{appendix:dataset}
The overall dataset collection and construction process is illustrated in Figure \ref{fig:data_collection_pipeline}
\begin{figure*}[t]
    \centering
    \includegraphics[width=\textwidth]{figures/data_collection_vqa_pipeline.pdf}
    \caption{Data collection and bilingual VQA benchmark construction pipeline. Plastic-waste images are collected from different locations in Dhaka, annotated using bounding boxes, organized into an annotated dataset, and converted into English--Bangla VQA pairs spanning recognition, visual description, and environmental reasoning.}
    \label{fig:data_collection_pipeline}
\end{figure*}
\subsection{Quality Control}
\label{subsec:quality_control}

The benchmark records image-quality and generation-status metadata. Of the 402 images, 385 are labeled \texttt{clear} and 17 \texttt{moderate}; none are retained as fully ambiguous. Seventeen images carry a low-image-quality warning, while one source has an empty label file and one has a missing label file.

Our construction is intentionally conservative but semi-automated. Category mappings and QA generation are not equivalent to independent expert annotation of every image. This distinction is essential: the benchmark should be interpreted as a controlled diagnostic dataset rather than a naturalistic human-authored VQA corpus.
\subsection{Task Hierarchy}
\label{subsec:task_hierarchy}

The benchmark contains exactly 804 QA pairs for each reasoning level.

\paragraph{Recognition (Easy).}
Questions ask which plastic-waste category is visually dominant. Answers are intentionally concise and category-centered.

\paragraph{Visual Description (Medium).}
Questions request observable properties such as general shape, flexibility, visible structure, openness, or deformation. Answers are designed to avoid claims that cannot be reliably inferred from visual evidence.

\paragraph{Environmental Reasoning (Hard).}
Questions ask why improper disposal may be environmentally problematic or what environmental concern is associated with the observed waste item. Reference answers use conservative concepts such as persistent litter, land and water pollution, movement into drains or waterways, wildlife exposure, and waste-management challenges.

Importantly, environmental answers do not claim exact degradation times, chemical toxicity, carbon footprint, contamination status, or direct medical consequences unless such evidence is externally established. 

\section{Additional Experimental Information}
\label{appendix:metrics}
\subsection{Zero-Shot Evaluation and Metrics}

All models are evaluated in a zero-shot setting without fine-tuning on the training or validation split. For the primary multimodal condition, each model receives the image together with the question in its original language and a short instruction encouraging direct, visually grounded, and scientifically conservative responses. A representative prompt is: ``Answer the question using the image. Use only visually supported information for recognition and description. For environmental questions, provide a concise scientifically conservative answer. Do not invent brands, exact chemical properties, locations, toxicity, or hidden objects.'' Temperature is set to 0 where supported, output length is capped to avoid unnecessary elaboration, and Bangla questions are processed directly without translation into English. For recognition, model outputs are normalized to the five target categories using a deterministic synonym mapping; for example, ``bottle,'' ``plastic bottle,'' and ``bottle-like plastic container'' are mapped to \texttt{plastic\_bottle}. Recognition performance is measured using micro accuracy,
\begin{equation}
\operatorname{Acc}
=
\frac{1}{N}
\sum_{i=1}^{N}
\mathbb{1}
\left[
\hat{y}_i = y_i
\right],
\label{eq:accuracy}
\end{equation}
where $y_i$ and $\hat{y}_i$ denote the reference and predicted categories, respectively, and macro-F1,
\begin{equation}
F1_{\mathrm{macro}}
=
\frac{1}{C}
\sum_{c=1}^{C}
F1_c,
\label{eq:macro_f1}
\end{equation}
where $C$ is the number of waste categories. Macro-F1 is reported because the benchmark is class-imbalanced and flexible-wrapper examples are dominant. For open-ended visual-description and environmental-reasoning responses, exact match is insufficient because semantically valid answers may differ in wording. We therefore combine semantic and lexical metrics. BERTScore-F1 is used as the primary automatic metric because it compares candidate and reference answers using contextual embeddings and is more tolerant of paraphrasing \citep{zhang2020bertscore}. BLEU-4 measures modified $n$-gram precision \citep{papineni2002bleu}, while ROUGE-L captures longest-common-subsequence overlap \citep{lin2004rouge}; both are reported as complementary lexical measures. Because English and Bangla differ in morphology, segmentation, and surface realization, lexical scores are interpreted cautiously and are not treated as perfectly language-neutral. For reproducibility, the experiments record the exact model or API snapshot, inference date, decoding parameters, package versions, and model configuration.
\subsection{Bilingual Consistency}
\label{subsec:bilingual_consistency}

Independent English and Bangla scores do not establish whether a model gives semantically equivalent answers to equivalent questions. We therefore compute a Bilingual Consistency Score (BCS). Since recognition, consistency is defined as category agreement across paired English and Bangla outputs. For open-ended answers, paired responses are represented in a multilingual sentence-embedding space using LaBSE, which was developed to provide language-agnostic sentence representations across more than 100 languages \citep{feng2022labse}. Given paired English and Bangla outputs $e_i$ and $b_i$, respectively, we compute

\begin{equation}
s_i =
\cos
\left(
\operatorname{LaBSE}(e_i),
\operatorname{LaBSE}(b_i)
\right).
\label{eq:labse_similarity}
\end{equation}

After validating a semantic threshold $\tau$ on a manually reviewed development subset, BCS is defined as

\begin{equation}
BCS
=
\frac{1}{N}
\sum_{i=1}^{N}
\mathbb{1}
\left[
s_i \geq \tau
\right].
\label{eq:bcs}
\end{equation}

The threshold $\tau$ is selected on a manually reviewed validation subset and then held fixed during test evaluation.

\subsection{Visual-Grounding Ablation}
\label{subsec:grounding_ablation}

Our central diagnostic removes the image while leaving the question unchanged.

The multimodal condition is

\begin{equation}
\mathcal{M}:
(I,Q)
\rightarrow
A,
\label{eq:multimodal_condition}
\end{equation}

whereas the language-only condition is

\begin{equation}
\mathcal{L}:
Q
\rightarrow
A.
\label{eq:language_condition}
\end{equation}

For an evaluation metric $S$, we define Visual Grounding Gain (VGG) as

\begin{equation}
VGG
=
S(\mathcal{M})
-
S(\mathcal{L}).
\label{eq:vgg}
\end{equation}

A large positive value indicates strong dependence on visual evidence. A small value suggests that a question can be answered largely through linguistic or domain priors.

This distinction is particularly important for environmental reasoning because general facts concerning plastic pollution may remain recoverable even when the model does not know whether the image contains a bottle, wrapper, cup, bag, or straw.


\subsection{Statistical Analysis}
\label{subsec:statistics}

Augmented variants originating from the same base image are statistically correlated. Treating all 62 test records as independent observations would therefore underestimate uncertainty. We use a cluster-bootstrap procedure based on \texttt{source\_base\_id}, rather than an ordinary QA-level bootstrap.
For each reported metric, base-image groups are sampled with replacement using bootstrap resampling. All augmented records and associated QA pairs belonging to a sampled base image are retained together. We report percentile-based 95\% confidence intervals.
Pairwise model comparisons use the same bootstrap samples across models, producing paired bootstrap estimates of performance differences.
\section{Additional Result}
\subsection{Error Analysis}
\label{subsec:error_analysis}

We manually audited a stratified subset of predictions covering models, languages, and question types. The annotation protocol distinguishes five failure categories.

\paragraph{Object misclassification.}
The dominant plastic item is assigned to an incorrect category.

\paragraph{Visual hallucination.}
The model invents a visible property, object, label, color, cap, opening, or condition that is not sufficiently supported by the image.

\paragraph{Generic environmental reasoning.}
The answer is environmentally plausible but nearly independent of the specific image or waste category.

\paragraph{Environmental hallucination.}
The model makes unsupported claims concerning toxicity, chemical composition, health consequences, or other scientific properties.

\paragraph{Bangla interpretation error.}
The model misinterprets the intent of the Bangla question or responds with semantically inappropriate content.

The observed qualitative error distribution is shown in Table~\ref{tab:error_distribution}.

\begin{table}[t]
\centering
\caption{Qualitative error distribution. }
\label{tab:error_distribution}
\begin{tabular}{lcc}
\toprule
\textbf{Error Type}
& \textbf{English}
& \textbf{Bangla} \\
\midrule
Object misclassification
& 6.5\%
& 9.7\% \\

Visual hallucination
& 5.4\%
& 8.1\% \\

Generic environmental reasoning
& 15.1\%
& \textbf{21.5\%} \\

Environmental hallucination
& 4.3\%
& 6.5\% \\

Bangla interpretation error
& --
& 7.5\% \\
\bottomrule
\end{tabular}
\end{table}

The observed error distribution identifies generic environmental reasoning as the dominant failure mode, particularly for Bangla. This behavior aligns with the visual-ablation analysis: models may possess broadly correct knowledge about plastic pollution while failing to condition that knowledge strongly on the precise visual evidence.

Each manually audited answer is additionally assessed using three scores in the range $[0,2]$:

\begin{equation}
H = C + R + G,
\label{eq:human_score}
\end{equation}

where $C$ denotes factual correctness, $R$ denotes relevance to the question, and $G$ denotes grounding in the supplied image. The maximum human-evaluation score is therefore 6. When multiple annotators are used, inter-annotator agreement is reported to quantify annotation consistency.

\section{Limitations and Threats to Validity}
\label{sec:limitations}
% 

\paragraph{Dataset size and diversity.}
The benchmark contains 136 underlying base images. The 402 image records include augmented variants, so the effective visual diversity is substantially smaller than the raw record count suggests. The benchmark should therefore be interpreted as a focused diagnostic resource rather than a comprehensive representation of global plastic pollution.

\paragraph{Category imbalance.}
Flexible wrappers account for 55.7\% of image records, whereas plastic bags represent only 6.0\%. Micro accuracy can therefore overstate performance on the dominant class. Macro-F1 and per-category evaluation reduce, but do not eliminate, this limitation.

\paragraph{Controlled rather than naturalistic language.}
Questions are generated through controlled templates and reference answers are drawn from a constrained answer bank. This improves consistency and enables paired bilingual comparisons but yields less linguistic diversity than human-authored VQA.

The current benchmark contains 2,412 unique question surface forms but only 100 distinct reference answer strings. Models may therefore benefit from repeated answer semantics. Future versions should incorporate independently authored references and greater paraphrastic diversity.

\paragraph{Annotation validity.}
Benchmark construction is semi-automated rather than independently expert-annotated at the full-dataset level. Before release as a high-stakes public benchmark, a stratified subset should be validated by human annotators for object category, question naturalness, Bangla quality, visual grounding, and scientific correctness.

\paragraph{Environmental scope.}
Environmental answers intentionally remain general and scientifically conservative. The benchmark does not evaluate polymer chemistry, life-cycle assessment, detailed ecological toxicity, local waste regulations, or contamination-specific risks. Models should therefore not be interpreted as environmental-science experts on the basis of this benchmark.

\paragraph{Model drift.}
Closed VLM APIs may change over time. Reproducibility requires fixed snapshots where possible and detailed recording of model identifiers, inference dates, prompts, decoding parameters, and provider-side changes.

\paragraph{Metric validity.}
BERTScore, BLEU, and ROUGE compare generated responses with textual references; none guarantees factual correctness. Similarly, LaBSE similarity may classify two fluent but factually incorrect responses as cross-lingually consistent. Human grounding evaluation remains necessary for high-confidence conclusions.

\paragraph{Dependence among augmented samples.}
Augmented images derived from the same base image are statistically dependent. Confidence intervals based on individual image records or QA pairs would therefore be artificially narrow. Our source-group cluster bootstrap explicitly addresses this dependence.


\end{document}